\ifdefined\niveldos\else
\documentclass[12pt,letterpaper]{report}
 
%Russian-specific packages
%--------------------------------------
\usepackage{bigints}
\usepackage[T2A]{fontenc}
\usepackage[utf8]{inputenc}
\usepackage[russian]{babel}
\usepackage[mathscr]{euscript}
\usepackage{amsmath,amsthm,amssymb,latexsym,amsfonts}
\usepackage{pythonhighlight}
\usepackage{lipsum}
\usepackage{soul}
\usepackage{amsmath,amssymb}
\usepackage{parskip}
\usepackage{graphicx}
\usepackage{mathtools}
\usepackage[usestackEOL]{stackengine} 
\usepackage{tocloft}
\usepackage{mdframed} % Required for boxed problems.
\usepackage{xcolor}
\usepackage{hyperref}
\usepackage{tikz}
\usepackage{thmtools}
\usepackage{amsthm}
\usepackage{mathtools}
\usepackage{pdfcomment}
\usepackage{blindtext}
\usepackage{todonotes}
\usepackage{import}
\usepackage{hyphsubst}
\usepackage{bookmark}
\usepackage{pgfplots}
\usepackage{tikzsymbols}
\usepackage{ tipa }
\usepackage{enumitem}
\usepackage{forloop}
\linespread{1.5}
\graphicspath{ {./images/} }
%%



\usepackage{amsmath}
\usepackage{graphicx}

\makeatletter
\newcommand{\ceq}{\mathbin{\mathpalette\make@smallcircled=}}
\newcommand{\ok}{\mathbin{\mathpalette\make@smallcircled\checkmark}}
\newcommand{\ola}{\mathbin{\mathpalette\make@circled\Longleftarrow}}
\newcommand{\ora}{\mathbin{\mathpalette\make@circled\Longrightarrow}}
\newcommand{\make@circled}[2]{%
  \ooalign{$\m@th#1\smallbigcirc{#1}$\cr\hidewidth$\m@th#1#2$\hidewidth\cr}%
}
\newcommand{\make@smallcircled}[2]{%
  \ooalign{$\m@th#1\smallerbigcirc{#1}$\cr\hidewidth$\m@th#1#2$\hidewidth\cr}%
}
\newcommand{\smallbigcirc}[1]{%
  \vcenter{\hbox{\scalebox{2.4}{$\m@th#1\bigcirc$}}}%
}
\newcommand{\smallerbigcirc}[1]{%
  \vcenter{\hbox{\scalebox{2}{$\m@th#1\bigcirc$}}}%
}
\makeatother

\makeatletter
\newcommand{\but}{\mathbin{\mathpalette\make@triangled!}}
\newcommand{\make@triangled}[2]{%
  \ooalign{$\m@th#1\smallbigtria{#1}$\cr\hidewidth$\m@th#1#2$\hidewidth\cr}%
}
\newcommand{\smallbigtria}[1]{%
  \vcenter{\hbox{\scalebox{2.4}{$\m@th#1\bigtriangleup$}}}%
}
\makeatother




\DeclarePairedDelimiter\abs{\lvert}{\rvert}%
\DeclarePairedDelimiter\norm{\lVert}{\rVert}%

% Swap the definition of \abs* and \norm*, so that \abs
% and \norm resizes the size of the brackets, and the 
% starred version does not.
\makeatletter
\let\oldabs\abs
\def\abs{\@ifstar{\oldabs}{\oldabs*}}
%
\let\oldnorm\norm
\def\norm{\@ifstar{\oldnorm}{\oldnorm*}}
\makeatother


\usepackage{tikz}
\usepackage{mathtools}
\newcommand{\myarrow}[1][-45]{%
  \mathrel{%
    \text{$
     \begin{tikzpicture}[baseline = -0.5ex]
       \node[inner sep=0pt,outer sep=0pt,rotate = #1] (a) at (0,0)  {$\xrightarrow{}$};
    \end{tikzpicture}
    $}%
  }%
}%

\newcommand{\inc}{
  \myarrow[45]
}

\newcommand{\dec}{
  \myarrow[-45]
}

\newcommand{\useg}[1] {
  \left\langle {#1} \right\rangle
}


{
  \bf
    \def\letus{%
    \mathord{\setbox0=\hbox{$\exists$}%
             \hbox{\kern 0.125\wd0%
                   \vbox to \ht0{%
                      \hrule width 0.75\wd0%
                      \vfill%
                      \hrule width 0.75\wd0}%
                   \vrule height \ht0%
                   \kern 0.125\wd0}%
           }%
}

\newenvironment{namedtheorem}[1]
{
\quad    
\begin{mdframed}[linewidth=0.6pt]
    \textbf{Теорема \textnormal{(#1)}:}
}
{\end{mdframed}}

\newenvironment{theorem}{
\quad
\begin{mdframed}[linewidth=0.6pt]
    
        \textbf{Теорема:}
}
{\end{mdframed}}

\newenvironment{properties}[1][,]
{
  \quad
  \textbf{Свойства:}
  \begin{enumerate}[#1]\raggedright
}
{
  \end{enumerate}
}

% \newlist{properties}{enumerate}{3}
% \setlist[properties]{before=\raggedright}
% \quad
% \begin{mdframed}[linewidth=0.6pt]
        % \textbf{Свойства:}
% }
% {\end{mdframed}}

\newenvironment{namedlemma}[1]
{
\quad    
\begin{mdframed}[linewidth=0.6pt]
    \textbf{Лемма \textnormal{(#1):}}
}
{\end{mdframed}}

\newenvironment{lemma}{
\quad
\begin{mdframed}[linewidth=0.6pt]
    
        \textbf{Лемма:}
}
{\end{mdframed}}


\newenvironment{example}{
    \textit{\underline{Пример:}}
}

\newtheorem*{remark}{Замечание}


\newenvironment{definition}{
    \quad
\begin{mdframed}[linewidth=0.6pt]
        \textbf{Определение:}
}
{\end{mdframed}}

% Definindo novas cores
\definecolor{verde}{rgb}{0.25,0.5,0.35}
\definecolor{jpurple}{rgb}{0.5,0,0.35}
\definecolor{darkgreen}{rgb}{0.0, 0.2, 0.13}
%\definecolor{oldmauve}{rgb}{0.4, 0.19, 0.28}
% Configurando layout para mostrar codigos Java
\usepackage{listings}

\newcommand{\estiloJava}{
\lstset{
    language=Java,
    basicstyle=\ttfamily\small,
    keywordstyle=\color{jpurple}\bfseries,
    stringstyle=\color{red},
    commentstyle=\color{verde},
    morecomment=[s][\color{blue}]{/**}{*/},
    extendedchars=true,
    showspaces=false,
    showstringspaces=false,
    numbers=left,
    numberstyle=\tiny,
    breaklines=true,
    backgroundcolor=\color{cyan!10},
    breakautoindent=true,
    captionpos=b,
    xleftmargin=0pt,
    tabsize=2
}}

\newcommand{\estiloR}{
  \lstset{ %
    language=R,                     % the language of the code
    basicstyle=\footnotesize,       % the size of the fonts that are used for the code
    numbers=left,                   % where to put the line-numbers
    numberstyle=\tiny\color{gray},  % the style that is used for the line-numbers
    stepnumber=1,                   % the step between two line-numbers. If it's 1, each line
                                    % will be numbered
    numbersep=5pt,                  % how far the line-numbers are from the code
    backgroundcolor=\color{white},  % choose the background color. You must add \usepackage{color}
    showspaces=false,               % show spaces adding particular underscores
    showstringspaces=false,         % underline spaces within strings
    showtabs=false,                 % show tabs within strings adding particular underscores
    frame=single,                   % adds a frame around the code
    rulecolor=\color{black},        % if not set, the frame-color may be changed on line-breaks within not-black text (e.g. commens (green here))
    tabsize=2,                      % sets default tabsize to 2 spaces
    captionpos=b,                   % sets the caption-position to bottom
    breaklines=true,                % sets automatic line breaking
    breakatwhitespace=false,        % sets if automatic breaks should only happen at whitespace
    title=\lstname,                 % show the filename of files included with \lstinputlisting;
                                    % also try caption instead of title
    keywordstyle=\color{blue},      % keyword style
    commentstyle=\color{darkgreen},   % comment style
    stringstyle=\color{red},      % string literal style
    escapeinside={\%*}{*)},         % if you want to add a comment within your code
    morekeywords={*,...}          % if you want to add more keywords to the set
}}

\newcommand{\vertiii}[1]{{\left\vert\kern-0.25ex\left\vert\kern-0.25ex\left\vert #1 
    \right\vert\kern-0.25ex\right\vert\kern-0.25ex\right\vert}}

\newcommand{\Z}{\mathbb{Z}}
\newcommand{\Q}{\mathbb{Q}}
\newcommand{\N}{\mathbb{N}}
\newcommand{\R}{\mathbb{R}}
\newcommand{\follow}{\textbf{\textit{Следствие:}}}
\newcommand{\notice}{\underline{\textit{Замечание }}}

\newcommand*\circled[1]{\tikz[baseline=(char.base)]{
            \node[shape=circle,draw,inner sep=2pt] (char) {#1};}}

\definecolor{linkcolor}{HTML}{47528f} % цвет ссылок
\definecolor{urlcolor}{HTML}{47528f} % цвет гиперссылок
\hypersetup{pdfstartview=FitH,  linkcolor=linkcolor,urlcolor=urlcolor, colorlinks=true}
\newcommand{\RomanNumeralCaps}[1]
  {\MakeUppercase{\romannumeral #1}}
\usepackage{titlesec}
\renewcommand{\thesection}{\arabic{section}}
\renewcommand{\listtheoremname}{Определения и теоремы}
\renewcommand\qedsymbol{$\blacksquare$}
\newcommand\oast{\stackMath\mathbin{\stackinset{c}{0ex}{c}{0ex}{\ast}{\bigcirc}}}
\makeatletter
\renewenvironment{proof}[1][\proofname]{%
   \par\pushQED{\qed}\normalfont%
   \topsep6\p@\@plus6\p@\relax
   \trivlist\item[\hskip\labelsep\bfseries#1\@addpunct{.}]%
   \ignorespaces\hfill\break$\square$
}{%
   \popQED\endtrivlist\@endpefalse
}
\makeatother
%--------------------------------------
\DeclareMathOperator{\Mr}{M_{\mathbb{R}}}
\addtolength{\oddsidemargin}{-.875in}
	\addtolength{\evensidemargin}{-.875in}
	\addtolength{\textwidth}{1.75in}

	\addtolength{\topmargin}{-.875in}
    \addtolength{\textheight}{1.75in}
%--------------------------------------
\def\calloutsym{%
  \ensurestackMath{%
  \scalebox{1.7}{\color{red}\stackunder[0pt]{\bigcirc}{\downarrow}}}%
}
\def\calloutsymup{%
  \ensurestackMath{%
  \scalebox{1.7}{\color{red}\stackon[0pt]{\bigcirc}{\uparrow}}}%
}
\newcommand\callouttext[1]{%
  \def\stacktype{S}\renewcommand\useanchorwidth{T}\stackText%
  \stackunder{\calloutsym}{\scriptsize\Longstack{#1}}\stackMath%
}
\newcommand\callout[3][2.5pt]{%
  \def\stacktype{L}\stackMath\stackunder[#1]{#2}{\callouttext{#3}}%
}
\newcommand\callouttextup[1]{%
  \def\stacktype{S}\renewcommand\useanchorwidth{T}\stackText%
  \stackon{\calloutsymup}{\scriptsize\Longstack{#1}}\stackMath%
}
\newcommand\calloutup[3][1.5pt]{%
  \def\stacktype{L}\stackMath\stackunder[#1]{#2}{\callouttextup{#3}}%
}
\begin{document}

\clearpage
%% temporary titles
% command to provide stretchy vertical space in proportion
\newcommand\nbvspace[1][3]{\vspace*{\stretch{#1}}}
% allow some slack to avoid under/overfull boxes
\newcommand\nbstretchyspace{\spaceskip0.5em plus 0.25em minus 0.25em}
% To improve spacing on titlepages
\newcommand{\nbtitlestretch}{\spaceskip0.6em}
\pagestyle{empty}
\begin{center}
\bfseries
\nbvspace[1]
\Huge
{\nbtitlestretch\huge
КОНСПЕКТ ЛЕКЦИЙ ПО МАТЕМАТИЧЕСКОМУ АНАЛИЗУ}

\nbvspace[1]
\normalsize

СПбГУ, МКН, СП\\
ЛЕКТОР: БАХАРЕВ ФЁДОР ЛЬВОВИЧ
\nbvspace[1]
\\
% \Large СОСТАВИТЕЛИ:\\[0.5em]
% \footnotesize 
% Андрей \href{https://github.com/K-dizzled}{K-dizzled} Козырев, 
% Никита  \href{https://github.com/muldrik}{muldrik} Митцев,
% Максим \href{https://github.com/maxmartynov08}{maxmartynov08} Мартынов,\\ 
% Семён \href{https://github.com/SmnTin}{SmnTin} Паненков,
% Марк \href{https://github.com/markprudnikov}{markprudnikov} Прудников,
% Мария \href{https://github.com/MashaK5}{MashaK5} Козловцева

\nbvspace[2]

\includegraphics[width=4.0in]{./images/logo.jpeg}
\nbvspace[3]
\normalsize

\large
2021-2023
\nbvspace[1]
\end{center}
\newpage
\pagestyle{plain}
\fi
%\maketitle
%%%%%%%%%%%%%%%%%%%%%%%%%%%%%%%%%%
%                                %
%                                %
%        Table of content        %
%                                %
%                                %
%%%%%%%%%%%%%%%%%%%%%%%%%%%%%%%%%%
\tableofcontents
%\listoftheorem-nons[ignore={lemma},show={conj, theorem-non}]
% To show Table of contents with
% definitions and lemmas
%----------------------------------------
% \listoftheorem-nons[ignoreall,show={lemma}]
% \listoftheorem-nons[ignoreall,show={conj}]
\newpage
\begin{normalsize}
\setcounter{section}{0}

% \newcounter{i}
% \forloop{i}{0}{\value{i} < 6}{
%   \include{sections/\arabic{i}}
% }

\section*{Лекция 1 (17.02.2022)}
\section{Интегральное исчисление функций}
\subsection{Первообразная и неопределенный интеграл}

\begin{definition}
    $f : (a, b) \rightarrow \R, F : (a, b) \rightarrow \R$, $F$ -- первообразная $f$ на $(a, b)$ если $F' = f$
\end{definition}


\begin{theorem}
    $f, F : (a, b) \rightarrow \R$, $F$ -- первообразная $f$ 
    \begin{enumerate}
    \item $\Rightarrow F + c$ - первообразная $f$
    \item $G : (a,b)$ - первообразная $f \Rightarrow F - G = const$ 
    \end{enumerate}
\end{theorem}

\begin{proof} \quad
    \begin{enumerate}
        \item Очевидно.
        \item $F' = f, G' = f \Rightarrow (F - G)' = 0 \Rightarrow (F - G) = const$(т.к. $g(a) - g(b) = g'(\theta)(a - b) = 0$ 
        по теореме Лагранжа)
    \end{enumerate}
\end{proof}

\begin{definition}
    $ \int f dx = $ {множество всех первообразных $f$} 
\end{definition}

\begin{properties}
    \item $\int f + g dx = \int f dx + \int g dx$
    
    \item $\int \lambda f dx = \lambda \int f dx, \lambda \neq 0$
        
    \item $\int \alpha f(x) + \beta g(x) dx = \alpha \int f(x) dx + \beta g(x) dx, \alpha^2 + \beta^2 \neq 0$
        
    \item (замена переменной в неопр. интеграле) \\
    $f : (a, b) \rightarrow \R, F$ - первообразная $f$, $\phi:(c, d) \rightarrow (a, b), 
    \phi$ - дифференцируема. \\
    $\Rightarrow \int f(\phi(x))\phi'(x) dx = F(\phi(x)) + c$
    \begin{remark}
        $\phi(x) = y, \phi'(x)dx = dy \Rightarrow \int f(y) dy = F(y) + c$
    \end{remark}

    \item(формула интегрирования по частям) \\
    \[f, g\text{ -- дифф на }(a, b) \Rightarrow \int f(x)g'(x)dx = fg - \int f'(x)g(x)dx\]
    \begin{proof}
        $(fg)' = fg' + f'g$
    \end{proof}

\end{properties}




\subsection{Площадь и псевдоплощадь}

\begin{definition}
    Определение площади: $S : 2 ^ {\R ^ 2} \rightarrow [0, +\infty)$
    \begin{enumerate}
        \item $S((a_1, b_1) \times (a_2, b_2)) = (b_1 - a_1) \cdot (b_2 - a_2)$
        \item $S(E_1 \sqcup E_2) = S(E_1) + S(E_2)$
    \end{enumerate}
\end{definition}

\begin{remark}
    Такая функция не существует.
\end{remark}

\begin{definition}(Псевдоплощадь) $\sigma : F \rightarrow [0, +\infty]$, где $F$ это ограниченные подмножества $R^2$
    \begin{enumerate}
        \item $\sigma(\langle a_1, b_1\rangle \times \langle a_2, b_2\rangle) = (b_1 - a_1) \cdot (b_2 - a_2)$

        \item $\sigma(E) = \sigma(E_-) + \sigma(E_+)$ если $E = E_- \sqcup E_+$ и $E_-, E_+$ получены разрезанием $E$ вертикальной или
        горизонтальной линией.

        \item $E_1 \supset E_2 \Rightarrow \sigma(E_1) \geq \sigma(E_2)$ 
    \end{enumerate}
    
\end{definition}

\begin{remark} 

    \begin{enumerate} 
        \item Не важно куда относить точки прямой, т.к. площадь прямой равна 0 \\
        \tikzset{every picture/.style={line width=0.5pt}} %set default line width to 0.75pt        
        \begin{tikzpicture}[x=0.75pt,y=0.75pt,yscale=-0.75,xscale=0.75]
        %uncomment if require: \path (0,300); %set diagram left start at 0, and has height of 300
        
        %Shape: Free Drawing [id:dp7628728695496989] 
        \draw  [color={rgb, 255:red, 0; green, 0; blue, 0 }  ][line width=3] [line join = round][line cap = round] (146.33,63.67) .. controls (136.73,63.67) and (135.51,59.39) .. (124.33,55.67) .. controls (113.51,52.06) and (98.23,50.72) .. (87.33,51.67) .. controls (79.49,52.35) and (72.38,64.62) .. (67.33,69.67) .. controls (47.68,89.32) and (34.42,122.83) .. (59.33,145.67) .. controls (89.45,173.27) and (145.78,137.25) .. (162.33,109.67) .. controls (165.06,105.12) and (163.56,98.3) .. (164.33,93.67) .. controls (164.49,92.74) and (166.24,92.61) .. (166.33,91.67) .. controls (166.93,85.36) and (167.63,78.87) .. (166.33,72.67) .. controls (164.68,64.75) and (143.8,59.67) .. (138.33,59.67) ;
        %Shape: Free Drawing [id:dp2860373614820697] 
        \draw  [color={rgb, 255:red, 0; green, 0; blue, 0 }  ][line width=3] [line join = round][line cap = round] (111.33,22.67) .. controls (100.18,56.13) and (103.91,94.37) .. (102.33,130.67) .. controls (102.08,136.41) and (95.91,157.36) .. (94.33,163.67) .. controls (94.1,164.58) and (92.38,164.73) .. (92.33,165.67) .. controls (91.7,178.32) and (92.33,191) .. (92.33,203.67) ;
        %Shape: Free Drawing [id:dp08947191931178611] 
        \draw  [color={rgb, 255:red, 0; green, 0; blue, 0 }  ][line width=3] [line join = round][line cap = round] (99.33,215.67) .. controls (107.27,215.67) and (119.57,211.43) .. (124.33,206.67) .. controls (126.86,204.14) and (129.54,190.09) .. (125.33,189.67) .. controls (101.17,187.25) and (115.33,216.77) .. (115.33,234.67) ;
        %Shape: Free Drawing [id:dp9334586901493984] 
        \draw  [color={rgb, 255:red, 0; green, 0; blue, 0 }  ][line width=3] [line join = round][line cap = round] (165.33,23.67) .. controls (165.33,32.82) and (163.33,42.46) .. (163.33,51.67) ;
        %Shape: Free Drawing [id:dp04768711693928851] 
        \draw  [color={rgb, 255:red, 0; green, 0; blue, 0 }  ][line width=3] [line join = round][line cap = round] (170.33,27.67) .. controls (173,27.67) and (175.95,26.47) .. (178.33,27.67) .. controls (179.53,28.26) and (178.33,30.33) .. (178.33,31.67) ;
        %Shape: Free Drawing [id:dp5274114829816234] 
        \draw  [color={rgb, 255:red, 0; green, 0; blue, 0 }  ][line width=3] [line join = round][line cap = round] (170.33,38.67) .. controls (171.8,38.67) and (177.33,38.74) .. (177.33,40.67) ;
        %Shape: Free Drawing [id:dp6570628135366894] 
        \draw  [color={rgb, 255:red, 0; green, 0; blue, 0 }  ][line width=3] [line join = round][line cap = round] (163.33,49.67) .. controls (169.95,49.67) and (176.03,51.67) .. (182.33,51.67) ;
        %Shape: Free Drawing [id:dp8681341271113268] 
        \draw  [color={rgb, 255:red, 0; green, 0; blue, 0 }  ][line width=3] [line join = round][line cap = round] (116.33,87.67) .. controls (116.33,94.92) and (118.33,102.33) .. (118.33,109.67) ;
        %Shape: Free Drawing [id:dp6050366900451063] 
        \draw  [color={rgb, 255:red, 0; green, 0; blue, 0 }  ][line width=3] [line join = round][line cap = round] (126.33,86.67) .. controls (133.01,86.67) and (139.66,85.67) .. (146.33,85.67) ;
        %Shape: Free Drawing [id:dp7624577504469063] 
        \draw  [color={rgb, 255:red, 0; green, 0; blue, 0 }  ][line width=3] [line join = round][line cap = round] (128.33,97.67) .. controls (132.39,97.67) and (140.33,95.61) .. (140.33,99.67) ;
        %Shape: Free Drawing [id:dp6175909634779915] 
        \draw  [color={rgb, 255:red, 0; green, 0; blue, 0 }  ][line width=3] [line join = round][line cap = round] (120.33,111.67) .. controls (128.09,111.67) and (131.2,108.67) .. (138.33,108.67) ;
        %Shape: Free Drawing [id:dp110598214336876] 
        \draw  [color={rgb, 255:red, 0; green, 0; blue, 0 }  ][line width=3] [line join = round][line cap = round] (143.33,107.67) .. controls (152.66,107.67) and (142.07,117.6) .. (142.33,117.67) .. controls (150.18,119.63) and (150.17,115.67) .. (156.33,115.67) ;
        %Shape: Free Drawing [id:dp36668561089277496] 
        \draw  [color={rgb, 255:red, 0; green, 0; blue, 0 }  ][line width=3] [line join = round][line cap = round] (62.33,91.67) .. controls (62.33,98.21) and (64.33,104.94) .. (64.33,111.67) ;
        %Shape: Free Drawing [id:dp2428588068124462] 
        \draw  [color={rgb, 255:red, 0; green, 0; blue, 0 }  ][line width=3] [line join = round][line cap = round] (64.33,93.67) .. controls (69.02,91.79) and (74.29,91.67) .. (79.33,91.67) ;
        %Shape: Free Drawing [id:dp21516400666861557] 
        \draw  [color={rgb, 255:red, 0; green, 0; blue, 0 }  ][line width=3] [line join = round][line cap = round] (69.33,99.67) .. controls (73.39,99.67) and (81.33,97.61) .. (81.33,101.67) ;
        %Shape: Free Drawing [id:dp3834606938145779] 
        \draw  [color={rgb, 255:red, 0; green, 0; blue, 0 }  ][line width=3] [line join = round][line cap = round] (70.33,109.67) .. controls (73.02,109.67) and (75.65,110.67) .. (78.33,110.67) ;
        %Shape: Free Drawing [id:dp5890643945584151] 
        \draw  [color={rgb, 255:red, 0; green, 0; blue, 0 }  ][line width=3] [line join = round][line cap = round] (82.33,112.67) .. controls (85.3,112.67) and (92,110.83) .. (94.33,109.67) .. controls (95.18,109.25) and (95.91,106.82) .. (96.33,107.67) .. controls (96.99,108.98) and (92.33,123.09) .. (92.33,124.67) ;
    
        \end{tikzpicture}

        \item Псевдоплощадь существует, но не единственна.
    \end{enumerate}
\end{remark}

\begin{example}
    $E \in F$ \\
    $P_j$ это прямоугольник со сторонами параллельными осям координат, но произвольных размеров.
    \begin{enumerate}
        \item $\sigma_1(E) = \inf \left\{ \sum \sigma_1(P_j) \mid E \subset \bigcup\limits_{j=1}^N P_j \right\}$
        \item $\sigma_2(E) = \inf \left\{ \sum \sigma_2(P_j) \mid E \subset \bigcup\limits_{j=1}^\infty P_j \right\}$
    \end{enumerate}
    \begin{remark}
        $\sigma_1(E) \geq \sigma_2(E)$ и если $K = ([0, 1] \cap \Q) \times [0, 1] \Rightarrow \sigma_1(K) = 1, \sigma_2(K) = 0$
    \end{remark}
\end{example}

\begin{theorem}
    $\sigma_1$ -- псевдоплощадь. 
\end{theorem}

\begin{proof}
    \begin{enumerate}
        \item Если прямоугольник покрыть конечным числом прямоугольников, то у покрытия сумма площадей не меньше площади прямоугольника.
        (Т.к. можно провести все вертикальные и горизонтальные линии и разбить на дизъюнктивное объединение)

        \item 
        \circled{$\geq$} $\sigma_1(E) \stackrel{?}{=} \sigma_1(E_-) + \sigma_1(E_+)$ \\
        $E \subset P_1 \cup \ldots \cup P_n$,  $P_j = P_j^- \cup P_j^+$,  $\sigma_1(P_j) = \sigma_1(P_j^-) + \sigma_1(P_j^+)$ \\
        Тогда $\sum\limits_{j=1}^N \sigma_1(P_j) = \sum\limits_{j=1}^n \sigma_1(P_j^-) + \sum\limits_{j=1}^n \sigma_1(P_j^+) \geq \sigma_1(E_-) + \sigma_1(E_+)$ \\
        $\Rightarrow \sigma_1(E) \geq \sigma_1(E_-) + \sigma_1(E_+)$ \\
        
        \circled{$\leq$} $P_1, \ldots, P_k$ - покрытие $E_-$ с точность $\epsilon$,  \\
        $P_{k + 1}, \ldots, P_n$ - покрытие $E_+$ с точность $\epsilon$ \\
        $\Rightarrow \sigma_1(E_-) + \epsilon \geq \sigma_1(P_1) + \ldots + \sigma_1(P_k)$ \\
        $\sigma_1(E_+) + \epsilon \geq \sigma_1(P_{k + 1}) + \ldots + \sigma_1(P_n)$ \\
        $\Rightarrow \sigma_1(E_-) + \sigma_1(E_+) + 2\epsilon \geq \sum\limits_{j = 1}^N \sigma_1(P_j) \geq \sigma_1(E)$

        \item Покрытие большего является покрытием меньшего.
        
         
    \end{enumerate}
\end{proof}

\begin{theorem}
    Псевдоплощадь инварианта относительно сдвигов.
\end{theorem}

\begin{proof}
    Покрытие также сдвинется.
\end{proof}

\begin{remark}
    Проверить то же самое для $\sigma_2$.
\end{remark}



 

\subsection{Определенный интеграл.}

Считаем, что зафиксирована псевдоплощадь $\sigma$.


\begin{definition}
    $a \in \R \Rightarrow a_+ = \max(a, 0), a_- = \max(-a, 0)$ \\
    $a_+ + a_- = |a|, a_+ - a_- = a$ \\
    Аналогично для функции $f$. \\

    $\Gamma_f = \left\{(x, f(x)) \in \R^2\right\}$ \\
    Для $f \geq 0 \  P_f = \left\{ (x, y) : y \in [0, f(x)]\right\}$ - подграфик f. \\
    Для $f \geq 0$ на $[a, b], P_f[a, b] = \left\{ (x, y) : x \in [a, b], y \in [0, f(x)]\right\}$ - подграфик f.
\end{definition}

\begin{remark}
    $f$ - непр $\Rightarrow f_+, f_-$ - непр.
\end{remark}

\begin{definition}
    $f$ - непр. на $[a, b]$, тогда \\
    \[\int_a^b f(x) dx = \sigma(P_{f_+} [a, b]) - \sigma(P_{f_-}[a, b])\]
\end{definition}

\begin{properties}
    \item $\int_a^a f \ dx = 0$
    \item $\int_a^b 0 \  dx = 0$
    \item $f \geq 0$ на $[a, b] \Rightarrow \int_a^b f \ dx \geq 0$
    \item $\int_a^b - f \ dx = -\int_a^b f \ dx$
    \item $\int_a^b c \ dx = c \cdot (b - a)$
    \item $f \geq 0$ на $[a, b] \wedge \int_a^b f \ dx = 0 \Rightarrow f = 0$ на $[a, b]$
    \begin{proof}
        Т.к. функция непрерывна, то если существует точка $f$ со значением не 0,
            то можно найти окрестность со значением $>$ 0 и там будет ненулевая площадь
    \end{proof}

    \item(Аддитивность интеграла) $c \in [a, b] \Rightarrow \int_a^b f dx = 
    \int_a^c f dx +\int_c^b f dx$
    
    \begin{proof}
        Следует из аддитивности псевдоплощади.
    \end{proof}

    \begin{remark}
        Соглашение: $\int_a^b f dx = - \int_b^a f dx \Rightarrow$ аддитивность верна и 
        для $c \not \in [a, b]$
    \end{remark}

    \item(Монотонность) $f \geq g$ на $[a ,b] \Rightarrow \int_a^b f dx \geq 
    \int_a^b g dx$
    
    \begin{proof}
        $f \geq g \Rightarrow f_+ \geq g_+ \wedge f_- \leq g_-$
    \end{proof}

    \item $(b - a) \min\limits_{[a, b]} f \leq \int_a^b f dx \leq (b - a) \max\limits_{[a, b]} f$
    \item(Теорема о среднем) $\exists c \in [a, b] : \int_a^b f dx = f(c) \cdot (b-a)$
    \begin{proof}
    \[\frac {\int_a^b f dx} {b - a} \in [\min f, \max f] \Rightarrow \exists c\]
    \end{proof}

    \item(Теорема Барроу) $\Phi(x) = \int_a^x f(t) dt, \Psi(x) = \int_x^b f(t) dt$ \\
    $\Phi$ - интеграл с переменным верхним пределом, $\Psi$ - нижним \\
    $\Phi' = f, \Psi' = -f$
    \begin{proof}
        $x_1, x_2$ 
        \[\frac {\Phi(x_1) - \Phi(x_2)} {x_1 - x_2}
        = \frac {\int_{x_1}^{x_2} f(t) dt} {x_1 - x_2}  
        \stackrel{x_1- fix}{\Rightarrow} \lim_{x_2 \to x_1} \frac {\Phi(x_1) - \Phi(x_2)} {x_1 - x_2}
        = \lim_{x_2 \to x_1} \frac {\int_{x_1}^{x_2} f(t) dt} {x_1 - x_2}
        = \lim_{x_2 \to x_1} f(\theta)\] \\
        где $\theta$ лежит между $x_1$ и $x_2$ \\
        $\Rightarrow \Phi(x_1) = f(x_1)$ \\
        $\Phi + \Psi = const \Rightarrow \Psi' = -f$
    \end{proof}

    \item(формула Ньютона-Лейбница) $F$ - первообразная $\Rightarrow$ \\
    $\int_a^b f(x) dx = F(b) - F(a) = F|_a^b$
    \begin{proof}
        \[ \Phi(x) = \int_a^x f \Rightarrow F = \Phi + c \Rightarrow F(b) - F(a)
        = \Phi(b) - \Phi(a) = \int_a^b f \]
    \end{proof}

    \item(Линейность интеграла) \[f, g \in C[a, b], \alpha, \beta \in R \Rightarrow
    \int_a^b (\alpha f + \beta g) dx = \alpha \int_a^b f dx + \beta \int_a^b g dx \] 
    \begin{proof}
        Следует из существования первообразной.
    \end{proof}
    \item $|\int_a^b f dx| \leq \int_a^b |f| dx$
    \begin{proof}
        \[ \left|\int_a^b f\right| = \left|\int_a^b f_+ - \int_a^b f_-\right| \leq
            \int_a^b f_+ + \int_a^b f_- = \int_a^b |f| \]
    \end{proof}
    
    \item(формула интегрирования по частям)
    \[f, x \in C^1[a, b] \Rightarrow
    \int_a^b f(x) g'(x) dx = f(x)g(x)\bigg|_a^b - \int_a^b f'(x)g(x) dx \]

    \begin{proof}
        Следует из формулы Ньютона-Лейбница и формулы интегрирования по частям для неопределенного интеграла.
    \end{proof}

    \item(Замена переменной) $ f:[a, b] \to \R, \phi [c, d] \to [a, b], \phi \in C^1[c, d]$ 
    \[\Rightarrow
    \forall p, q \in [c, d] \int_p^q f(\phi(t)) \phi'(t) dt
    = \int_{\phi(t)}^{\phi(q)} f(x) dx\]

    \begin{example}
        \begin{enumerate}
            \item \[ \int_0^1 e^{\sin(x)} \cos(x) \ dx = [\sin(x) = y, dy = \cos(x) dx] 
            = \int_{\sin(0)}^{sin(1)} e^y \ dy = \ldots \]
            \item \[ \int_0^1 \sqrt{1 - x^2} \ dx = [x = \sin(y), dx 
            = \cos(y) dy] = \int_0^{\arcsin(1)} \sqrt{1 - \sin^2(y)} \cos(y) \ dy 
            = \int_0^{\frac \pi 2} \cos^2(y) \ dy \] 
            \[ = \int_0^{\frac \pi 2} \frac {1 + \cos^2(y)} 2 \ dy
            = \int_0^{\frac \pi 2} \frac 1 2 \ dy = \frac \pi 4\]
        \end{enumerate}
    \end{example}
\end{properties}


\section*{Лекция 2 (24.02.2022)}

\begin{example}
 \[
     \begin{gathered}
         W_n = \int_0^{\pi / 2}{(\sin{x})^n}\ dx = \int_0^{\pi / 2}{(\cos{x})^n}\ dx 
          \text{ (т.к. можно сделать замену $y = \pi / 2 - x$)}
         \\
         W_n = \int_0^{\pi / 2}\cos^n{x}\ dx = \int_0^{\pi / 2}\cos^{n - 1}{x}\ d(\sin{x}) =
         \cos^{n - 1}{x} \cdot \sin{x} \bigg|_{0}^{\pi/2} - \int_{0}^{\pi / 2}\sin{x}\ d(\cos^{n - 1}{x}) = \\
         = 0 \text{ (при $n \geqslant 2$)} - \int_0^{\pi / 2} \sin{x} \cdot (n - 1) \cos^{n - 2}{x}(-\sin{x})\ dx =
         (n - 1)\int_{0}^{\pi / 2}(1 - \cos^2{x}) \cos^{n - 2}{x}\ dx =\\ 
         = (n - 1)(W_{n - 2} - W_{n})
     \end{gathered}
\]

Итого получаем рекуррентную формулу:
\[
    W_n = \frac{n - 1}{n} \cdot W_{n - 2}, \, W_0 = \pi / 2, \, W_1 = 1
\]
Тогда итоговая формула:
\[
    \begin{gathered}
        W_{2k} = \frac{2k - 1}{2k} \cdot \frac{2k - 3}{2k - 2} \cdot \ldots \cdot \frac{1}{2} \cdot W_0 =
        \frac{(2k - 1)!!}{(2k)!!} W_0 = \frac{(2k - 1)!!}{(2k)!!} \pi / 2
        \\\\
        W_{2k + 1} = \frac{2k}{2k + 1} \cdot \frac{2k - 2}{2k - 1} \cdot \ldots \cdot \frac{2}{3} \cdot W_1 =
        \frac{(2k)!!}{(2k + 1)!!} W_1 = \frac{(2k)!!}{(2k + 1)!!} 
    \end{gathered}
\]

Также, стоит отметить, что $W_n \geqslant W_{n + 1} \geqslant W_{n + 2} \geqslant \dots$ т.к.
$\sin^n{x} \leqslant \sin^{n - 1}{x}$, то есть интегрируем меньшую функцию, на одном и том же промежутке.
Также, заметим, что это на самом деле есть неплохое приближение числа $\pi$.\\

\underline{\textit{Формула Валлиса}}
\[
    \begin{gathered}
        W_{2k + 2} \leqslant W_{2k + 1} \leqslant W_{2k}\\
        \Downarrow\\
        \frac{\pi}{2} \cdot \frac{(2k + 1)!!}{(2k + 2)!!} \leqslant \frac{(2k)!!}{(2k + 1)!!} \leqslant
        \frac{\pi}{2} \cdot \frac{(2k - 1)!!}{(2k)!!}\\
        \Downarrow\\
        \frac{\pi}{2} \cdot \frac{2k + 1}{2k + 2} \leqslant \frac{((2k)!!)^2}{(2k + 1)((2k - 1)!!) ^ 2} \leqslant \frac{\pi}{2}\\
        \Downarrow\\
        \lim_{k \to \infty} \frac{(2k!!)^2}{(2k + 1)((2k - 1)!!) ^ 2} = \pi / 2\\
        \Downarrow\\
        \lim_{k \to \infty} \frac{2k!!}{\sqrt{2k + 1}(2k - 1)!!} = \sqrt{\pi / 2}
    \end{gathered}
\]
\follow 
\[
    C_{2k}^k \sim \sqrt{\frac{2}{\pi}} \cdot \frac{4 ^ k}{\sqrt{2k + 1}}
\]
\begin{proof}
     Можно заметить, что
    \[
        (2k)!! = 2 ^ k \cdot k!,\qquad
        (2k + 1)!! = \frac{(2k + 1)!}{(2k)!!} = \frac{(2k + 1)!}{2^k \cdot k!}
    \]
    Откуда можно получить следующее равенство:
    \[
        \begin{gathered}
            \frac{(2k)!!}{(2k - 1)!!} = \frac{(2^k \cdot k!)^2}{(2k)!} = \frac{4 ^ k \cdot (k!)^2}{(2k)!}
            = \frac{4 ^ k}{C_{2k}^k}, \text{последнее равенство верно так как:}\\
            \sum_{k = 0}^{2n} C_{2n}^{k} = 4 ^ n, \,
            C_{2n}^n = \frac{(2n)!}{(n!)^2}\\
        \end{gathered}
    \]
    Тогда по формуле Валлиса:
    \[
        \frac{2k!!}{\sqrt{2k + 1}(2k - 1)!!} \to \sqrt{\frac{\pi}{2}}
        \Leftrightarrow 
        \frac{1}{\sqrt{2k + 1}} \cdot \frac{4 ^ k}{C_{2k}^k} \to \sqrt{\frac{\pi}{2}}
        \Leftrightarrow 
        C_{2k}^k \sim \sqrt{\frac{2}{\pi}} \cdot \frac{4 ^ k}{\sqrt{2k + 1}}
    \]
\end{proof}
\end{example}
\newpage

\section{Формула Тейлора с остатком в интегральной форме}

\quad
\[
    \begin{gathered}
        f(x) = T_{n, x_0} f(x) + R_{n, x_0} f(x),\\
        \text{где } T_{n, x_0} f(x) = f(x_0) + f'(x_0)(x - x_0) + \dots + \frac{f^{(n)}(x_0)}{n!}(x - x_0)^n
    \end{gathered}
\]

\begin{namedtheorem}{остаток в интегральной форме}
    
    Если \, $f \in C^{n + 1} \langle a, b \rangle$,  \,\, $x, x_0 \in \langle a, b \rangle$, то
    \[
        R_{n, x_0}f(x) = \frac{1}{n!} \int_{x_0}^x f^{(n + 1)}(t)(x-t)^n\ dt
    \]
\end{namedtheorem}

\begin{proof}

    \textbf{\large База:} 
    \[
        n = 0, \text{ то } f(x) = f(x_0) + \int_{x_0}^{x}f'(t)\ dt  \text{ --- формула Ньютона-Лейбница}
    \]

    \textbf{\large Переход:} 
    \[
        \begin{gathered}
        R_{n, x_0}f(x) = \frac{1}{n!} \int_{x_0}^{x}f^{(n+1)}(t)(x - t)^n\ dt =
        \frac{1}{n!} \int_{x_0}^x f^{(n + 1)}(t)\ d\left(-\frac{(x - t) ^ {n + 1}}{n + 1}\right) =\\\\
        = - \frac{1}{n!} \frac{(x - t) ^ {n + 1}}{n + 1} f^{(n + 1)}(t) \bigg|_{x_0}^{x} +
        \frac{1}{n!} \int_{x_0}^{x} \frac{(x - t)^{n + 1}}{n + 1}\ d\left(f^{(n + 1)}(t)\right) =\\\\
        = \frac{(x - x_0) ^ {n + 1}}{(n + 1)!} f^{n + 1}(x_0) + R_{n + 1, x_0} f(x)
        \end{gathered}
    \]


\end{proof}


\begin{namedlemma}{чуть более хитрая теорема о среднем}
    
    $f, g \in C[a, b], g \geqslant 0 \Rightarrow \exists c \in (a, b) : \int_a^b f(x) g(x) dx = f(c) \int_a^b g(x) dx$
\end{namedlemma}

\newpage 
\begin{proof}
    \[
        \begin{gathered}
            m = \min_{[a, b]} f, \,\,\, M = \max_{[a, b]} f\\
            m\ g(x) \leqslant f(x) g(x) \leqslant M\ g(x)\\
            m \int_a^b g(x)\ dx \leqslant \int_a^b f(x) g(x)\ dx \leqslant M \int_a^b g(x)\ dx\\
            \frac{\int_a^b f(x) g(x)\ dx}{\int_a^b g(x)\ dx} \in [m, M] \Rightarrow 
            \text{равно значению $f$ в некоторой точке $c\in [a,b]$}
        \end{gathered}
    \]
\end{proof}
\begin{follow}
        
    Берём остаток в интегральной форме:
    \[
        R_{n, x_0} f(x) = \frac{1}{n!} \int_{x_0}^x f^{(n + 1)}(t) (x - t) ^ n\ dt
    \]
    Заметим, что $(x - t)^n$ имеет одинаковый знак на $t\in [x_0, x]$. Тогда по только что доказанной лемме $\exists \theta$
    , что остаток можно переписать в следующей форме.
    \[
        \frac{1}{n!} f^{(n + 1)} (\theta) \int_{x_0}^x (x - t) ^ n dt = \frac{f^{(n + 1)}(\theta)}{(n + 1)!}  \cdot (x - x_0) ^ {(n + 1)}
    \]
    Таким образом, остаток в форме Лагранжа без проблем выводится из остатка в интегральной форме.
\end{follow}

\begin{remark}
    Некоторые студенты на экзамене стали жертвой следующего:
    \[
        \ln{(1+x)} = x - \frac{x ^ 2}{2} + \frac{x ^ 3}{3} - ... + R_n
    \]
    \[
        R_n = \frac{n!}{(n + 1)!} \cdot \frac{(-1) ^ {n + 1}}{(1 + \theta) ^ n} \cdot x^{n + 1} = \frac{(-1) ^ {n + 1}}{n + 1} \cdot \frac{x ^ {n + 1}}{(1 + \theta) ^ n}
    \]
    \quad 
    При каких $x$,\, $R_n \to 0$?
    При $x \in (0, 1)$ утверждение очевидно, но что делать при $x \in (-1, 0)$?
    \[
        R_n = \frac{1}{n!} \int_0^x ... \text{ To be continued (Федор Львович думал, что так можно, но не получилось :) )}
    \]
\end{remark}
\newpage 

\subsection{Приближенное вычисление интеграла}

\begin{definition}
    $[a, b], \tau $ --- дробление отрезка $[a, b]$ $\tau = \{ x_0, x_1, \ldots, x_n \}$, \\ где $a = x_0 < x_1 < \ldots < x_n = b$ \quad $|\tau| = \max\limits_{k = 0, \ldots, n - 1}{|x_{k + 1} - x_k|}$ - мелкость (ранг) дробления

    $\theta = \{ t_1, \ldots, t_n\}$ -- оснащение дробления $\tau$, где $t_j \in [x_{j - 1}, x_j]$

    $(\tau, \theta) $ --- оснащенное дробление

    $f \in C[a, b]$ интегральная сумма $S_{\tau, \theta} (f) = \sum\limits_{j = 0}^{n - 1} f(t_{j + 1}) (x_{j + 1} - x_j)$
\end{definition}

\quad

% Pattern Info
 
\tikzset{
pattern size/.store in=\mcSize, 
pattern size = 5pt,
pattern thickness/.store in=\mcThickness, 
pattern thickness = 0.3pt,
pattern radius/.store in=\mcRadius, 
pattern radius = 1pt}
\makeatletter
\pgfutil@ifundefined{pgf@pattern@name@_fhsu0lo8n}{
\pgfdeclarepatternformonly[\mcThickness,\mcSize]{_fhsu0lo8n}
{\pgfqpoint{0pt}{0pt}}
{\pgfpoint{\mcSize+\mcThickness}{\mcSize+\mcThickness}}
{\pgfpoint{\mcSize}{\mcSize}}
{
\pgfsetcolor{\tikz@pattern@color}
\pgfsetlinewidth{\mcThickness}
\pgfpathmoveto{\pgfqpoint{0pt}{0pt}}
\pgfpathlineto{\pgfpoint{\mcSize+\mcThickness}{\mcSize+\mcThickness}}
\pgfusepath{stroke}
}}
\makeatother
\tikzset{every picture/.style={line width=0.75pt}} %set default line width to 0.75pt        
\usetikzlibrary{patterns}
\begin{tikzpicture}[x=0.75pt,y=0.75pt,yscale=-1,xscale=1]
%uncomment if require: \path (0,300); %set diagram left start at 0, and has height of 300

%Shape: Axis 2D [id:dp02029949369434092] 
\draw  (132,233.4) -- (534.5,233.4)(172.25,66) -- (172.25,252) (527.5,228.4) -- (534.5,233.4) -- (527.5,238.4) (167.25,73) -- (172.25,66) -- (177.25,73)  ;
%Straight Lines [id:da7710677350311679] 
\draw    (223,222) -- (223,242) ;
%Straight Lines [id:da6636852601941825] 
\draw    (290,223) -- (290,243) ;
%Straight Lines [id:da500250663548284] 
\draw    (267,224) -- (267,244) ;
%Straight Lines [id:da377281395464541] 
\draw    (243,222) -- (243,242) ;
%Straight Lines [id:da8386061762887231] 
\draw    (311,223) -- (311,243) ;
%Straight Lines [id:da7815801182452056] 
\draw    (333,224) -- (333,244) ;
%Straight Lines [id:da8616890669141692] 
\draw    (345,225) -- (345,245) ;
%Shape: Free Drawing [id:dp5560591048641171] 
\draw  [color={rgb, 255:red, 0; green, 0; blue, 0 }  ,draw opacity=1 ][line width=0.75] [line join = round][line cap = round] (222,195) .. controls (243.45,162.83) and (256.46,175.18) .. (301,172) .. controls (305.12,171.71) and (312.5,172.5) .. (317,171) .. controls (327.51,167.5) and (340.44,159.22) .. (343,149) .. controls (343.86,145.54) and (344,120.8) .. (344,128) ;
%Straight Lines [id:da9019023552875143] 
\draw [color={rgb, 255:red, 208; green, 2; blue, 27 }  ,draw opacity=1 ] [dash pattern={on 4.5pt off 4.5pt}]  (223,193) -- (223,222) ;
%Straight Lines [id:da3323747111843628] 
\draw [color={rgb, 255:red, 208; green, 2; blue, 27 }  ,draw opacity=1 ] [dash pattern={on 4.5pt off 4.5pt}]  (243,175) -- (243,192) -- (243,222) ;
%Straight Lines [id:da04332372509475069] 
\draw [color={rgb, 255:red, 208; green, 2; blue, 27 }  ,draw opacity=1 ] [dash pattern={on 4.5pt off 4.5pt}]  (267,174) -- (267,224) ;
%Straight Lines [id:da18598253754944805] 
\draw [color={rgb, 255:red, 208; green, 2; blue, 27 }  ,draw opacity=1 ] [dash pattern={on 4.5pt off 4.5pt}]  (290,173) -- (290,223) ;
%Straight Lines [id:da9229530047708664] 
\draw [color={rgb, 255:red, 208; green, 2; blue, 27 }  ,draw opacity=1 ] [dash pattern={on 4.5pt off 4.5pt}]  (311,173) -- (311,223) ;
%Straight Lines [id:da27194960638393784] 
\draw [color={rgb, 255:red, 208; green, 2; blue, 27 }  ,draw opacity=1 ] [dash pattern={on 4.5pt off 4.5pt}]  (332,163) -- (333,224) ;
%Straight Lines [id:da8974921338138321] 
\draw [color={rgb, 255:red, 208; green, 2; blue, 27 }  ,draw opacity=1 ] [dash pattern={on 4.5pt off 4.5pt}]  (344,125) -- (345,225) ;
%Straight Lines [id:da25565245177699003] 
\draw [color={rgb, 255:red, 74; green, 144; blue, 226 }  ,draw opacity=1 ]   (223,193) -- (243,193) ;
%Straight Lines [id:da5135245657261043] 
\draw [color={rgb, 255:red, 74; green, 144; blue, 226 }  ,draw opacity=1 ]   (243,175) -- (267,174) ;
%Straight Lines [id:da40438608893718364] 
\draw [color={rgb, 255:red, 74; green, 144; blue, 226 }  ,draw opacity=1 ]   (267,174) -- (290,173) ;
%Straight Lines [id:da8267126008581839] 
\draw [color={rgb, 255:red, 74; green, 144; blue, 226 }  ,draw opacity=1 ]   (290,173) -- (311,173) ;
%Straight Lines [id:da6583698282228275] 
\draw [color={rgb, 255:red, 74; green, 144; blue, 226 }  ,draw opacity=1 ]   (311,173) -- (332,173) ;
%Straight Lines [id:da6130580258949248] 
\draw [color={rgb, 255:red, 74; green, 144; blue, 226 }  ,draw opacity=1 ]   (332,163) -- (346,163) ;
%Straight Lines [id:da9281438699801965] 
\draw    (276,227) -- (276,244) ;
%Shape: Rectangle [id:dp8370752348674357] 
\draw  [pattern=_fhsu0lo8n,pattern size=8.100000000000001pt,pattern thickness=0.75pt,pattern radius=0pt, pattern color={rgb, 255:red, 126; green, 211; blue, 33}] (267,174) -- (290,174) -- (290,233) -- (267,233) -- cycle ;

% Text Node
\draw (252,245.4) node [anchor=north west][inner sep=0.75pt]    {$x_{j}$};
% Text Node
\draw (291,246.4) node [anchor=north west][inner sep=0.75pt]    {$x_{j+1}$};
% Text Node
\draw (272,248.4) node [anchor=north west][inner sep=0.75pt]    {$t_{j}$};


\end{tikzpicture}

\begin{theorem}
    $f \in C[a, b] \Rightarrow S_{\tau, \theta}(f) \underset{|\tau| \to 0} {\to} \int_a^b f(x) dx$, 
    \\ \\
    т.е. $\forall \varepsilon > 0 \,\, \exists \delta > 0 : \forall \tau, \theta : | \tau | < \delta \Rightarrow \abs{ \int_a^b f dx - S_{\tau, \theta}(f)} < \varepsilon$
\end{theorem}

\quad

\begin{definition}
    $f$ функция на $[a, b]$, существование предела $S_{\tau, \theta}(f)$ при $\abs{\tau} \to 0 $ означает, что $f$ интегрируема по Риману и $\lim = \int_a^{b} f dx$
\end{definition}

\quad 

\underline{\textit{Упражнение:}}  Доказать, что если у функции есть один разрыв первого рода, то она интегрируема по Риману


\begin{proof}

    По теореме Кантора $\forall \varepsilon > 0 \,\, \exists \, \delta > 0 : \forall x, y \in [a, b] : \abs{x - y} < \delta \Rightarrow \abs{f(x) - f(y)} < \frac{\varepsilon}{b - a}$

    Берем дробление $\tau : \abs{\tau}  < \delta $
    \[
        \begin{gathered}
            \abs{S_{\tau, \theta} - \int\limits_{a}^{b}{f(x)\ dx}}=
            \abs{\sum_{j = 1}^n f(t_j) (x_j - x_{j - 1}) - \sum_{j = 1}^n \int_{x_{j - 1}}^{x_j} f(x)\ dx} =\\
            =\text{(по теореме о среднем)}=\abs{\sum_{j = 1}^n f(t_j) (x_j - x_{j - 1}) - \sum_{j = 1}^n f(c_j) \cdot (x_j - x_{j - 1})} = \\
            = \abs{
                \sum_{j = 1} ^ n (f(t_j) - f(c_j)) \cdot (x_j - x_{j - 1})
            } < \frac{\varepsilon}{b - a} \sum_{j = 1} ^ n (x_j - x_{j - 1}) = \varepsilon
        \end{gathered}
    \]
\end{proof}


\underline{\textit{Задача:}}  $f$ имеет разрыв в точке $c$ первого рода и $f \in C[a, b), f \in C(c,b], \exists \lim\limits_{x \to c-}{f}, \exists \lim\limits_{x \to c+}{f}$, тогда $S_{\tau, \theta} f \to \int_a^c f + \int_c^b f$

\quad 

\begin{example}
    \begin{enumerate}
    \item 
    \[
        \begin{gathered}
            \int_0^{a} e^x dx = \,? , \,\,\, \tau = \left\{ 0, \frac{a}{n}, \frac{2a}{n}, \ldots, \frac{na}{n} \right\}, \,\,\, \theta = \left\{ 0, \frac{a}{n}, \ldots , \frac{n - 1}{n} a \right\}\\
            S_{\tau, \theta} f = \sum_{j = 1}^n f \left( \frac{a_{j - 1}}{n} \right)(x_j - x_{j - 1}) 
            = \frac{a}{n} \cdot \left( 1 + e ^ {\frac{1}{n}a} + \ldots + e ^{\frac{n - 1}{n} a} \right) 
            = \frac{a}{n} \cdot \frac{e ^ {\frac{na}{n}} - 1}{e ^ {\frac{a}{n}} - 1} \to e^a - 1
        \end{gathered}
    \]
    \item
    \[
        \begin{gathered}
            1 ^ p + 2 ^ p + \ldots + n ^ p = n ^ {p + 1} \cdot \frac{1}{n} \left( \left( \frac{1}{n} \right) ^ p + \ldots + \left( \frac{n}{n} \right) ^ p \right) \, \ceq\\
            \tau = \left\{ 0, \frac{1}{n}, \ldots , \frac{n}{n} \right\}, \theta = \left\{\frac{1}{n}, \ldots , \frac{n}{n}  \right\}, f(x) = x ^ p\\
            \ceq \, n ^ {p + 1} S_{\tau, \theta} f = \left[\frac{1 ^ p + \ldots + n ^ p}{(n + 1) ^ p} \to \int_0^1 x ^ p dx\right]
            = \frac{n^{p + 1}}{p + 1} \Rightarrow 1 ^ p + \ldots + n ^ p \sim \frac{n ^ {p + 1}}{p + 1}
        \end{gathered}
    \]

    \item интеграл Пуассона
    \[
        \int_0^\pi \log{(1 - 2a \cos{x} + a ^ 2)} dx
    \]

    \[
        f = \log{(1 - 2a \cos{x} + a ^ 2)},\, [0, \pi],\, \tau = \left\{0, \frac{\pi}{n}, \ldots , \frac{\pi n}{n} \right\}, \theta = \left\{0, \frac{\pi}{n}, \ldots , \frac{\pi (n - 1)}{n} \right\}
    \]

    \[
        S_{\tau, \theta} f = \sum_{j = 1}^n \frac{\pi}{n} \cdot \log\left(1 - 2a \cos{\frac{\pi (j - 1)}{n} + a ^ 2}\right) = \frac{\pi}{n} \log\left(\prod_{j = 1}^n \left(1 - 2a \cos{\frac{\pi (j - 1)}{n} + a ^ 2}\right)\right)
    \]

    Заметим, что $a ^ 2 - 2a \cos{\theta} + 1 = (a - e ^ {-i \theta})(a - e ^ {i \theta})$:
    \[
        \begin{gathered}
            \prod_{j = 1}^n \left(1 - 2a \cos{\frac{\pi (j - 1)}{n} + a ^ 2}\right)=
            \prod_{j = 1}^n \left(a - e ^ {\frac{-\pi (j - 1)}{n} i }\right)\left(a - e ^ {\frac{\pi (j - 1)}{n} i }\right)=\\
            =\left[(z^{2n} - 1) = \prod\limits_{j=0}^{2n}{z - e^{i\frac{j \pi}{n}}}\right]=
            \frac{(a ^ {2n} - 1) \cdot (a - 1)}{(a + 1)} 
        \end{gathered}
    \]
    Тогда, использовав полученное равенство в исходном уравнение, получим, что 
    \[
        \frac{\pi}{n} \log\left(\prod_{j = 1}^n \left(1 - 2a \cos{\frac{\pi (j - 1)}{n} + a ^ 2}\right)\right)=
        \frac{\pi}{n} \cdot \log{\frac{(a ^ {2n} - 1)(a - 1)}{a + 1}} \underset{n \to \infty}{\to}
        \begin{cases}
            0 & \abs{a} < 1 \\
            2 \pi \log{\abs{a}} & \abs{a} > 1 \\
            ? & \abs{a} = 1
        \end{cases}
    \]
    \end{enumerate}
    
\end{example}


\begin{remark}
    $\int_a^b f(x)\ dx \approx S_{\tau, \theta}(f)$
    И можно оценить погрешность следующим образом:
     \[
        \abs{\int\limits_{a}^{b}{f(x)\ dx} - S_{\tau, \theta}(f)} \leqslant \max{\abs{f'}} \cdot \abs{\tau} (b - a) 
    \] 
\end{remark}
\begin{proof}
    Пусть $t_j, c_j \in [x_j, x_{j + 1}]$, где $t_j$~--- точки выбранного оснащения $\theta$, а для $c_j$ выполнено следующее:
    \[
        \int\limits_{x_j}^{x_{j + 1}}{f(x)\ dx} = f(c) (x_{j + 1} - x_j)
    \]
    Тогда, пользуясь выкладками из прошлой теоремы, можно утверждать, что:
     \[
        \abs{\int\limits_{a}^{b}{f(x)\ dx} - S_{\tau, \theta}(f)} =
        \sum\limits_{j = 1}^{n} \abs{f(t_j) - f(c_j)}(x_{j + 1} - x_j) 
    \] 
    Тогда, если мы знаем $\max{\abs{f'}}$, то, пользуясь тем, что $t_j, c_j \in [x_j, x_{j + 1}] \Rightarrow \abs{t_j - c_j} \le \abs{\tau}$,
    можно дать следующую оценку:
    \[
        \begin{gathered}
            \sum\limits_{j = 1}^{n} \abs{f(t_j) - f(c_j)}(x_{j + 1} - x_j) \leqslant 
            \sum\limits_{j = 1}^{n} \abs{t_j - c_j} \max{\abs{f'}}(x_{j + 1} - x_j) \leqslant 
            \max{\abs{f'}} \cdot \abs{\tau} (b - a) 
        \end{gathered}
    \]
\end{proof} 
\quad 

\begin{remark}
    Более точная формула для приближенного вычисления интеграла
\end{remark}


% Pattern Info
 
\tikzset{
pattern size/.store in=\mcSize, 
pattern size = 5pt,
pattern thickness/.store in=\mcThickness, 
pattern thickness = 0.3pt,
pattern radius/.store in=\mcRadius, 
pattern radius = 1pt}
\makeatletter
\pgfutil@ifundefined{pgf@pattern@name@_hmcbrwike}{
\pgfdeclarepatternformonly[\mcThickness,\mcSize]{_hmcbrwike}
{\pgfqpoint{0pt}{0pt}}
{\pgfpoint{\mcSize+\mcThickness}{\mcSize+\mcThickness}}
{\pgfpoint{\mcSize}{\mcSize}}
{
\pgfsetcolor{\tikz@pattern@color}
\pgfsetlinewidth{\mcThickness}
\pgfpathmoveto{\pgfqpoint{0pt}{0pt}}
\pgfpathlineto{\pgfpoint{\mcSize+\mcThickness}{\mcSize+\mcThickness}}
\pgfusepath{stroke}
}}
\makeatother
\tikzset{every picture/.style={line width=0.75pt}} %set default line width to 0.75pt        

\begin{tikzpicture}[x=0.75pt,y=0.75pt,yscale=-1,xscale=1]
%uncomment if require: \path (0,300); %set diagram left start at 0, and has height of 300

%Shape: Axis 2D [id:dp02029949369434092] 
\draw  (132,233.4) -- (534.5,233.4)(172.25,66) -- (172.25,252) (527.5,228.4) -- (534.5,233.4) -- (527.5,238.4) (167.25,73) -- (172.25,66) -- (177.25,73)  ;
%Straight Lines [id:da7710677350311679] 
\draw    (223,222) -- (223,242) ;
%Straight Lines [id:da6636852601941825] 
\draw    (290,223) -- (290,243) ;
%Straight Lines [id:da500250663548284] 
\draw    (267,224) -- (267,244) ;
%Straight Lines [id:da377281395464541] 
\draw    (243,222) -- (243,242) ;
%Straight Lines [id:da8386061762887231] 
\draw    (311,223) -- (311,243) ;
%Straight Lines [id:da7815801182452056] 
\draw    (333,224) -- (333,244) ;
%Straight Lines [id:da8616890669141692] 
\draw    (345,225) -- (345,245) ;
%Shape: Free Drawing [id:dp5560591048641171] 
\draw  [color={rgb, 255:red, 0; green, 0; blue, 0 }  ,draw opacity=1 ][line width=0.75] [line join = round][line cap = round] (222,195) .. controls (243.45,162.83) and (256.46,175.18) .. (301,172) .. controls (305.12,171.71) and (312.5,172.5) .. (317,171) .. controls (327.51,167.5) and (340.44,159.22) .. (343,149) .. controls (343.86,145.54) and (344,120.8) .. (344,128) ;
%Straight Lines [id:da9019023552875143] 
\draw [color={rgb, 255:red, 208; green, 2; blue, 27 }  ,draw opacity=1 ] [dash pattern={on 4.5pt off 4.5pt}]  (223,193) -- (223,222) ;
%Straight Lines [id:da3323747111843628] 
\draw [color={rgb, 255:red, 208; green, 2; blue, 27 }  ,draw opacity=1 ] [dash pattern={on 4.5pt off 4.5pt}]  (243,175) -- (243,192) -- (243,222) ;
%Straight Lines [id:da04332372509475069] 
\draw [color={rgb, 255:red, 208; green, 2; blue, 27 }  ,draw opacity=1 ] [dash pattern={on 4.5pt off 4.5pt}]  (267,174) -- (267,224) ;
%Straight Lines [id:da18598253754944805] 
\draw [color={rgb, 255:red, 208; green, 2; blue, 27 }  ,draw opacity=1 ] [dash pattern={on 4.5pt off 4.5pt}]  (290,173) -- (290,223) ;
%Straight Lines [id:da9229530047708664] 
\draw [color={rgb, 255:red, 208; green, 2; blue, 27 }  ,draw opacity=1 ] [dash pattern={on 4.5pt off 4.5pt}]  (311,173) -- (311,223) ;
%Straight Lines [id:da27194960638393784] 
\draw [color={rgb, 255:red, 208; green, 2; blue, 27 }  ,draw opacity=1 ] [dash pattern={on 4.5pt off 4.5pt}]  (332,163) -- (333,224) ;
%Straight Lines [id:da8974921338138321] 
\draw [color={rgb, 255:red, 208; green, 2; blue, 27 }  ,draw opacity=1 ] [dash pattern={on 4.5pt off 4.5pt}]  (344,125) -- (345,225) ;
%Straight Lines [id:da25565245177699003] 
\draw [color={rgb, 255:red, 74; green, 144; blue, 226 }  ,draw opacity=1 ]   (223,193) -- (243,175) ;
%Straight Lines [id:da5135245657261043] 
\draw [color={rgb, 255:red, 74; green, 144; blue, 226 }  ,draw opacity=1 ]   (243,175) -- (267,174) ;
%Straight Lines [id:da40438608893718364] 
\draw [color={rgb, 255:red, 74; green, 144; blue, 226 }  ,draw opacity=1 ]   (267,174) -- (290,173) ;
%Straight Lines [id:da8267126008581839] 
\draw [color={rgb, 255:red, 74; green, 144; blue, 226 }  ,draw opacity=1 ]   (290,173) -- (311,173) ;
%Straight Lines [id:da6583698282228275] 
\draw [color={rgb, 255:red, 74; green, 144; blue, 226 }  ,draw opacity=1 ]   (311,173) -- (332,163) ;
%Straight Lines [id:da6130580258949248] 
\draw [color={rgb, 255:red, 74; green, 144; blue, 226 }  ,draw opacity=1 ]   (332,163) -- (344,125) ;
%Straight Lines [id:da9281438699801965] 
\draw    (276,227) -- (276,244) ;
%Shape: Rectangle [id:dp8370752348674357] 
\draw  [pattern=_hmcbrwike,pattern size=8.100000000000001pt,pattern thickness=0.75pt,pattern radius=0pt, pattern color={rgb, 255:red, 126; green, 211; blue, 33}] (267,174) -- (290,174) -- (290,233) -- (267,233) -- cycle ;

% Text Node
\draw (252,245.4) node [anchor=north west][inner sep=0.75pt]    {$x_{j}$};
% Text Node
\draw (291,246.4) node [anchor=north west][inner sep=0.75pt]    {$x_{j+1}$};
% Text Node
\draw (272,248.4) node [anchor=north west][inner sep=0.75pt]    {$t_{j}$};


\end{tikzpicture}

Приближаем трапециями (формула трапеций)

\[
    \int_a^b f dx  \approx \sum (x_{j + 1} - x_j) \cdot \frac{f(x_j) + f(x_{j + 1})}{2}
\]

\begin{namedtheorem}{о погрешности в формуле трапеций}
$f \in C^2[a, b]$

\[
    \abs{\int_a^b f - \sum (x_{j + 1} - x_j) \cdot \frac{f(x_j) + f(x_{j + 1})}{2}} \leqslant \frac{1}{8} \abs{\tau} ^ 2 \int_a^b \abs{f''(x)}dx
\]
\end{namedtheorem}


\section*{Лекция 5 (17.03.2022)}
\section{Кривые и пути}
\subsection{Вводная в кривые}
\begin{definition}
    Путь в $\R^m, \gamma : [a, b] \mapsto \R^m$~--- непрерывное отображение из отрезка в $\R^m$.
\end{definition}
\begin{definition}
    Носителем пути или кривой назовём следующее множество точек: $\gamma\left([a, b]\right)\in \R^m$
\end{definition}
\begin{definition}
    Путь замкнут если $\gamma(a) = \gamma(b)$. Иными словами путь замкнут, если начало пути совпадает с концом.
\end{definition}
\begin{definition}
    Если $\gamma$~--- инъекция, то такой путь называется
    простым.
\end{definition}
\begin{remark}
    Можно определять путь покоординатно:
    $\gamma(t) = (\gamma_1(t),\dots, \gamma_m(t))$ 
    Где $\gamma_i$~--- координатные функции.
\end{remark}
\begin{definition}
    Назовём два пути $\gamma_1\colon [a,b]\mapsto \R^m$ и
    $\gamma_2\colon [c,d]\mapsto \R^m$ эквивалентными,\\
    если $\exists \text{строго возрастающая биекция }
    \varphi\colon [a,b]\mapsto[c, d]$, так что 
    $\gamma_2\circ \varphi = \gamma_1$.
\end{definition}
\begin{remark}
    Определение эквивалентности путей является корректным отношением эквивалентности для пространства путей в $\R^m$.

    Можно воспринимать это следующим образом.
    Кривая в $\R_m$~--- класс эквивалентности путей. Представители~--- параметризация отрезком.
\end{remark}
\begin{example}
    У носителя пути есть различные параметризации:
    $(\cos x, \sin x), (x, \sqrt{1 - x^2})$(полуокружность)
\end{example}
\begin{remark}
    Зная параметризацию можно определить носитель кривой,
    начало и конец, но хочется понимать ещё про гладкость.
\end{remark}
\begin{definition}
    $\gamma$ ~--- гладкий путь, если $\gamma_i\in C_1[a,b]$.\\
\end{definition}
\begin{definition}
    Гладкая кривая~--- кривая, у которой $\exists$ гладкий путь.
\end{definition}
\begin{remark}
    Нельзя утверждать, что у гладкой кривой есть явная касательная
    в каждой её точке.
\end{remark}
\subsection{Длина кривой}
По факту, хотим подробить нашу кривую на маленькие отрезочки
и засуммировать их, давайте подумаем как это сделать аккуратно.\\
\quad
\begin{definition}
    $\gamma\colon [a, b]\mapsto \R^m$~--- путь.
    $\theta = \{t_0,\dots, t_n\}, a = t_0 < t_1 < \dots < t_n = b$~---
    разбиение $[a, b]$.
    \[
        l_\theta(\gamma) = \sum\limits_{j = 1}^{n} \abs{\gamma(t_j) - \gamma(t_{j - 1})};\\
    \]
    Тогда длиной назовём  $l(\gamma) = \sup_\theta l_\theta(\gamma)$
\end{definition}
\begin{remark}
    Заметим, что длина есть у любого пути, так как мы не вводили
    дополнительные требования на путь.
\end{remark}
\begin{remark}
    Посмотрим на $(x, x\sin \frac{1}{x})$
    Вопрос: Путь бесконечной длины?
\end{remark}
\begin{namedtheorem}{Независимость длины от параметризации}
    $\gamma_1 \sim \gamma_2\Rightarrow l(\gamma_1) = l(\gamma_2)$
\end{namedtheorem}
\begin{proof}
    $\theta$~--- дробление $[a,b]$,
    $\varphi(\theta)$~--- дробление $[c,d]$.
    $\gamma_2\circ\varphi=\gamma_1$.
    \[
        \sum\limits_{j = 1}^{n}\abs{\gamma_1(t_j) - \gamma_1(t_{j-1})}=\\
        \sum\limits_{j=1}^{n}\abs{\gamma_2(\varphi(t_j))-\gamma_2(\varphi(t_{j-1}))}
        \Rightarrow
        l_{\theta}(\gamma_1) = l_{\varphi(\theta)}(\gamma_2)
    \]
\end{proof}
\begin{namedtheorem}{Аддитивность пути}
    $\gamma\colon [a,b]\mapsto\R^m, c\in(a,b)$.
    Пусть $\gamma_{-}=\gamma\big|_{[a,b]}, \gamma_{+} = \gamma\big|_{[b,c]}$.
    Требуется доказать, что $l(\gamma)=l(\gamma_-) + l(\gamma_+)$
\end{namedtheorem}
\begin{proof}
    % TODO
    Докажем неравенство в две стороны.
    \begin{itemize}
        \item
        Пусть $\theta = \{x_1,\dots, x_n\}$~--- произвольное дробление $[a,b]$.
        Выберем за $k$ следующее число: $x_k \in [a,c], x_{k + 1}\in [c,b]$.
        Разделим наше дробления на два:
        $\theta_+ = (\theta \cap [a,c]) \cup \{c\} = \{x_1,\dots,x_k, c\},
        \theta_- = \{c\} \cup (\theta \cap [c, b]) = \{c, x_{k+1},\dots, x_n\}$.
        \[
            \begin{gathered}
                l_{\theta}(\gamma) =
                \sum\limits_{i=1}^{n}\norm{\gamma(x_i) - \gamma(x_{i + 1})}=
                \sum\limits_{i=1}^{k-1}\norm{\gamma(x_i) - \gamma(x_{i + 1})} + \norm{x_k - x_{k+1}} +
                \sum\limits_{i=k+1}^{n-1}\norm{\gamma(x_i) - \gamma(x_{i + 1})}\le\\
                \le\sum\limits_{i=1}^{k-1}\norm{\gamma(x_i) - \gamma(x_{i + 1})}
                + \norm{x_k - c} + \norm{c - x_{k + 1}} +
                \sum\limits_{i=k+1}^{n-1}\norm{\gamma(x_i) - \gamma(x_{i + 1})}=
                l_{\theta_-}(\gamma) + l_{\theta_+}(\gamma)
            \end{gathered}
        \] 
        Так как было взята произвольная $\theta$, то из этого следует и оценка на $l(\gamma) = \sup_\theta(\dots)$.
        \[
            l(\gamma) \le l(\gamma_+) + l(\gamma_-)
        \] 
        \item
        Докажем в обратную сторону.
        Пусть $\theta_-, \theta_+$~--- дробления $[a,c], [c, b]$ соответственно.
        Тогда определим $\theta = \theta_- \cup \theta_+$. Отсюда, по определению длины следует неравенство:
         \[
             l(\gamma) = \sup_{\tau}(l_{\tau}(\gamma)) \ge l_{\theta}(\gamma) = l_{\theta_-}(\gamma) + l_{\theta_+}(\gamma)
        \] 
        Так как были взяты произвольные $\theta_-, \theta_+$, то можно утверждать, что:
         \[
             l(\gamma) \ge l(\gamma_+) + l(\gamma_-)
        \] 
    \end{itemize}
\end{proof}
\begin{namedtheorem}{O длине гладкого пути}
    $\gamma\colon[a,b]\mapsto\R^m, \gamma_j\in C^1[a,b] \Rightarrow 
    l(\gamma)=\bigintss_{a}^{b}\abs{\gamma'(t)}dt$, где
    $\abs{\gamma'(t)} = \sqrt{\sum\limits_{j=1}^{m}{\abs{\gamma'_j(t)}}^2}$
\end{namedtheorem}
\begin{proof}
    $\theta$ ~--- дробление $[a,b]$,  
    $\theta\colon a = t_0<\dots<t_n=b$ 
    $$l_{\theta}(\gamma)=\sum\limits_{j=1}^{m}\abs{\gamma(t_j) - \gamma(t_{j-1})}$$
     \[
         \abs{\gamma(t_j)-\gamma(t_{j-1})}=
         \sqrt{\sum\limits_{k=1}^{m}\abs{\gamma_k(t_j)-\gamma_k(t_{j-1})}^2}=
         \sqrt{\sum\limits_{k=1}^{m}{\abs{\gamma'_k(?)}^2(t_j-t_{j-1})^2}}=
         (t_j - t_{j-1})\sqrt{\sum\limits_{k=1}^{m}{\abs{\gamma'_k(?)}^2}}
     .\] 
     Пусть $m_{j,k}=\min\limits_{[t_{j-1},t_j]} \abs{\gamma_k'}$,
     $M_{j,k}=\max\limits_{[t_{j-1},t_j]}\abs{\gamma'_k}$.
     Тогда можно продолжить  следующим образом:
     $$(t_j - t_{j-1})\sqrt{\sum\limits_{k=1}^{m}{m_{j,k}^2}}\leq
     \abs{\gamma(t_j) - \gamma(t_{j-1})}\leq
     (t_j - t_{j-1})\sqrt{\sum\limits_{k=1}^{m}{M_{j,k}^2}}$$
     С другой стороны:
     $$(t_j - t_{j-1})\sqrt{\sum\limits_{k=1}^{m}{m_{j,k}^2}}\leq
     \int\limits_{t_{j-1}}^{t_j}{\abs{\gamma'(t)}dt}\leq
     (t_j - t_{j-1})\sqrt{\sum\limits_{k=1}^{m}{M_{j,k}^2}}$$
     Суммируем по $j$:
     \[
         \sum\limits_{}^{}{(t_j-t_{j-1})\sqrt{?}}\leq
         \begin{cases}
             l_{\theta}(\gamma)\\
             \int\limits_{a}^{b}{\abs{\gamma'(t)}dt}
         \end{cases}\leq
         \sum\limits_{}^{}{(t_j-t_{j-1})\sqrt{?}}
     .\] 
     Заметим, что обе части стремятся к одному и тоже при
     $\abs{\theta}\rightarrow 0$, так как в этом случае $m_{j,k}\approx M_{j,k}$.

     Тогда для  $\varepsilon > 0$ по теореме кантора $\delta > 0\colon\\
     \forall s,t\in [a,b]\colon \abs{s - t} < \delta, 
     \forall k = 1\dots m, \abs{\gamma'_k(s)-\gamma'_k(t)} < \varepsilon$ 
     Тогда $\abs{M_{j,k}-m_{j,k}}<\varepsilon$ если $t_j - t_{j-1}<\delta$.
     \[
         \sqrt{\sum\limits_{k=1}^{m}{M^2_{j,k}}}-
         \sqrt{\sum\limits_{k=1}^{m}{m^2_{j,k}}} \leq
         \sqrt{\sum\limits_{k=1}^{m}{(M_{j,k} - m_{j,k})^2}}<
         \sqrt{m}\varepsilon
     .\] 
     Тогда $$\abs{\text{п.ч - л.ч}} <
     \sum\limits_{k=1}^{m}{(t_j-t_{j-1})\sqrt{m}\varepsilon}=
     (b - a)\sqrt{m}\varepsilon$$
\end{proof}
\begin{follow}
    \begin{enumerate}
        \item (Длина графика функции)\\
            Пусть $f\colon[a,b]\mapsto \R, f\in C^1[a,b], 
            G_f=\{(x,f(x)) \mid x\in [a,b]\}\subseteq \R^2$.
            Тогда длина графика функции равна:
            $l(G_f)=\bigintss\limits_{a}^{b}{\sqrt{1 + (f'(x))^2}dx}$
        \item
            Есть функция 
            $r\colon [\alpha, \beta]\mapsto \R_+, r\in C^1[\alpha, \beta]$
            Тогда можно запараметризовать эту кривую как:
            $(r(t)\cos t, r(t)\sin t)$, обозначим за координаты
            $x(t), y(t)$ соответственно.

            Тогда:
             $x'(t) = r'(t)\cos t - r(t)\sin t; 
             y'(t) = r'(t)\sin t + r(t)\cos t$

             Значит длину кривой можно вычислить следующим образом:
             \[
                 \begin{gathered}
                     l = 
                     \int\limits_{\alpha}^{\beta}{\sqrt{(x'(t))^2 + (y'(t))^2}}=\\
                     =\int\limits_{\alpha}^{\beta}{\sqrt{(\cos^2 t + \sin^2 t)(r'(t))^2 + (\cos^2 t + \sin^2 t)r^2(t)}\ dt}=\\
                     =\int\limits_{\alpha}^{\beta}{\sqrt{r'(t)^2 + r(t)^2}\ dt}
                 \end{gathered}
             \]
    \end{enumerate}
\end{follow}
\section{Линейные и полилинейные отображения}
\begin{remark}
    Прежде чем начать, подумаем что нам вообще надо.
    Производную можно представлять как линейное приближение
    функции в какой-либо точке.
    $f\colon X\mapsto Y$, где $X,Y$~--- линейные, нормированные,
    полные, пространства над одним полем
    скаляров($\R$ или $\mathbb{C}$).
    $f(x) = f(x_0) + A(x-x_0) + o(x - x_0)$.
    Объясним почему наложены те или иные условия.
    В одномерье $A$ было числом. В общем случае $A$ должно
    быть линейным отображением, поэтому нам хочется чтобы поле
    скаляров у $X,Y$ было одним и тем же.
    Более того, хочется иметь возможность переходить к пределу,
    а когда мы переходим к пределу естественно хотеть иметь
    корректное расстояние, поэтому нам нужна нормированность и полнота.
\end{remark}
\subsection{Основные свойства и определения}
\begin{definition}
    $X, Y$ ~--- линейные пространства над одним полем скаляров.
    $U\colon X\mapsto Y$ ~--- линейное, если
    \begin{enumerate}
        \item (Аддитивность)\\
            $U(x_1 + x_2)=U(x_1)+U(x_2)$
        \item (Однородность)\\
            $U(\lambda x) = \lambda U(x)$
    \end{enumerate}
\end{definition}
\begin{definition}
    Пусть $X_1, X_2, \dots, X_n, Y$ ~--- пространства над одним 
    полем скаляров.
    Тогда $U\colon X_1\times\dots\times X_n\mapsto Y$~--- полилинейное,
    если оно линейно по каждому из аргументов.
\end{definition}
\begin{remark}
    Часто скобочки опускаются: $U(x) = Ux$
\end{remark}
\begin{example}
    \begin{enumerate}
        \item
            $X = C[-1, 1], \delta\colon X\mapsto \R, \delta(f)=f(0)$. 
            Тогда $\delta$~--- линейное отображение.
        \item
        $X = C[a,b], Y = \R$. Тогда 
        $Uf = \bigintss\limits_{a}^{b}{fdx}$ ~---
        линейное отображение.
        \item
            $X = C[a,b], Y = C[a,b]$. Тогда
           $(Uf)(x) = \bigintss\limits_{a}^{x}{f(t)dt}$ 
        \item
            $X = C^1[a,b], Y = C[a,b]$. Тогда 
            $(Df)(x) = f'(x), D\colon X\mapsto Y$ тоже
            линейное отображение.
        \item
            $X_1 = X_2 = \dots X_n = \R = Y$. Тогда
            $U(x_1,\dots, x_n) = x_1\cdot \dots \cdot x_n$~---
            полилинейное отображение.
        \item
           $X_1 = \R^m, X_2 = \R^m, Y = \R$. Тогда
           $U(X_1, X_2) = (X_1, X_2)$, где $(X_1, X_2)$~---
           скалярное произведение, линейное по первой координате.
        \item
            $X_1 = \R^3, X_2 = \R^3, Y = \R^3$. Тогда
            $U(x_1, x_2) = x_1\times x_2 = [x_1, x_2]$, где
            $[x_1,x_2]$~--- линейное отображение,
            полилинейное отображение.
        \item
            $X_1 = \dots = X_m = \R^m$. Тогда 
            $U(x_1,\dots, x_m) = \det(x_1, \dots, x_m)$ 
            полилинейное отображение.
    \end{enumerate}
\end{example}
\begin{namedtheorem}{О эквивалентных условиях непрерывности}
    $U\colon X\mapsto Y$ ~--- линейное, $X,Y$ ~--- линейные,
    нормированные отображения (далее лно)
    над одним полем скаляров.
    Тогда следующие утверждения эквивалентны:
    \begin{enumerate}
        \item$U$ ~--- непрерывно
        \item $U$ ~--- непрерывно в 0
        \item $\exists C\colon \forall x\in X, \abs{\abs{Ux}}_Y\leq C\abs{\abs{x}}_X$
    \end{enumerate}
\end{namedtheorem}
\begin{proof}
    \begin{enumerate}
        \item $(1)\Rightarrow (2)$~--- очевидно
        \item  $(2)\Rightarrow (3)$
             $\forall \varepsilon > 0, \exists \delta > 0\colon
             \forall x\colon \norm{x}\leq\delta, \norm{Ux}\leq \varepsilon$.
             $x \rightarrow \overline{x} = x\frac{\delta}{\norm{x}}
             \Rightarrow \norm{U\overline{x}}\leq \varepsilon$ 
             % TODO недоделано
        \item $(3)\Rightarrow(1)$
             % $\|U(x_1) - U(x_2)\| \leq C\|x_1 - x_2\|$~---
             липшецевость  $\Rightarrow$ непрерывное.
    \end{enumerate}
\end{proof}
\begin{theorem}
    $U\colon X_1\times\dots \times X_n\mapsto Y$~--- полилинейное.
    Тогда следующие утверждения эквивалентны:
    \begin{enumerate}
        \item
            $U$ ~--- непрерывно
        \item 
            $U$ ~--- непрерывно в 0
        \item
            $\exists C\colon \norm{U(X_1,\dots, X_n)} \leq
            C\norm{x_1}\cdot\ldots\cdot\norm{x_n}$
    \end{enumerate}
\end{theorem}
\begin{remark}
    % TODO
\end{remark}
\begin{definition}
    $U: X\rightarrow Y$ ~--- лно.
    $\norm{U} = \inf\{C \mid \forall x\in X \norm{Ux}\leq C\norm{x}\}$
\end{definition}
\begin{remark}
$\inf$ достигается то есть: $\forall x\in X, \norm{Ux}\leq
    \norm{U}\cdot \norm{x}$
\end{remark}
\begin{example}
    $U\colon C[a,b]\mapsto C[a,b]$, причём
    $(Uf)(x) = \bigintss\limits_{a}^{x}{f(t)dt}$.
    Хотим оценить.
    % TODO
\end{example}


\import{sections/}{new.tex}
$ $

\end{normalsize}
\end{document} 
